\documentclass[11pt]{article}
\usepackage[margin=1in]{geometry}
\usepackage{amsmath,amssymb,amsthm}
\usepackage{mathtools}
\usepackage{hyperref}
\hypersetup{colorlinks=true,linkcolor=blue,citecolor=blue,urlcolor=blue}

\newtheorem{theorem}{Theorem}
\newtheorem{lemma}{Lemma}
\newtheorem{proposition}{Proposition}
\newtheorem{corollary}{Corollary}
\theoremstyle{definition}
\newtheorem{definition}{Definition}
\newtheorem{remark}{Remark}

\newcommand{\R}{\mathbb{R}}
\newcommand{\dd}[2]{\frac{\partial #1}{\partial #2}}

\title{Reciprocal Covariance Pencils and Equal-Specific-Damping\\ Symmetry in Two-Mode Inertial Ornstein--Uhlenbeck Systems}
\author{Elias De Jes\'us\\[2pt] \small Independent Researcher\\[2pt] \small ORCID: \href{https://orcid.org/0009-0007-0190-9143}{0009-0007-0190-9143}}
\date{September 2026}

\begin{document}
\maketitle

\begin{abstract}
We study a two-mode inertial (position--momentum) Ornstein--Uhlenbeck
system driven by two independent thermal baths, one coupled to each
oscillator's momentum coordinate. Writing $C_1,C_2$ for the stationary
covariance matrices produced by each bath acting alone, the joint
stationary covariance is the affine pencil $\Sigma(T_1,T_2)=T_1C_1+T_2C_2$,
and for the bath-scale ratio $s=T_1/T_2$ the quartic
$P(s)=\det(sC_1+C_2)$ governs whether exchanging the two bath scales
($s\mapsto 1/s$) induces a reciprocal symmetry of the stationary
covariance determinant. We show that, on the genuinely coupled interior
of parameter space, $P(s)$ is palindromic if and only if the two
oscillators share equal specific damping $\delta_i=\gamma_i/m_i$,
equivalently $\gamma_1/m_1=\gamma_2/m_2$. This is obtained from a
stronger, oriented result: the log-determinant comparator
$F=\log(\det C_2/\det C_1)$ is strictly monotone in $\delta_1$ at fixed
$\omega_1,\omega_2,k,\delta_2$, so that
$\operatorname{sign}F=-\operatorname{sign}(\delta_1-\delta_2)$
throughout the admissible domain, with a single zero crossing exactly at
$\delta_1=\delta_2$. At that balance locus we further establish the exact
identity $\partial F/\partial\delta_1|_{\delta_1=\delta_2=\delta}=-4/\delta$,
independent of the oscillator frequencies and coupling strength. We give
the exact positivity certificate underlying global monotonicity, record
an equivalent reciprocal-generalized-eigenvalue characterization of
palindromicity via $C_1^{-1}C_2$, and reduce the palindromic quartic to
an explicit quadratic in $w=s+s^{-1}$ on the balance locus. Several of the
polynomial identities used below were established by exact computer
algebra (symbolic simplification to zero and zero-remainder polynomial
division, never floating-point sampling); this is disclosed explicitly
where it occurs, together with the accompanying verification scripts.
\end{abstract}

\section{Introduction}

A first-order (overdamped) pair of coupled Ornstein--Uhlenbeck processes
driven by two independent baths produces a stationary covariance whose
determinant, viewed as a function of the bath-scale ratio $s$, is
quadratic in $s$; on the relevant reciprocal locus it reduces to an
affine relation in $w=s+s^{-1}$. That first-order theory is developed
separately~\cite{dejesus2026heterogeneity} and formalized in this
repository's companion Lean development.

Here we ask the analogous question for the \emph{inertial} extension of
that system: two damped harmonic oscillators, each carrying its own
position and momentum coordinate, linearly coupled through a shared
spring constant, each driven by an independent thermal bath acting on its
own momentum coordinate. The stationary covariance is now a $4\times4$
matrix, and the bath-pencil determinant $P(s)=\det(sC_1+C_2)$ is a genuine
quartic in $s$. We ask exactly when this quartic is palindromic --
equivalently, when its four roots (necessarily real and positive, in a
sense made precise below) close under $s\mapsto1/s$.

The answer is that the criterion is \emph{equal specific damping}
$\delta_1=\delta_2$, $\delta_i:=\gamma_i/m_i$ -- not equal raw damping
constants $\gamma_i$ and not equal masses $m_i$ individually. We prove
this as an unconditional \emph{if and only if} statement on the
genuinely-coupled interior of parameter space (Theorem~\ref{thm:main}),
via a stronger oriented result (Theorem~\ref{thm:monotone}) giving the
sign of the log-determinant comparator $F=\log(\det C_2/\det C_1)$
everywhere, not merely at the palindromic locus. A targeted literature
search did not locate an exact prior formulation of this criterion, the
oriented comparator, or the balance-slope identity of
Proposition~\ref{prop:slope}; the underlying machinery -- stationary
covariances of linear systems via the Lyapunov equation
\cite{kailath1980}, palindromic-polynomial algebra, and structured
matrix-pencil eigenvalue theory \cite{mackey2006} -- is classical or a
known tool, cited where it is used below. We report the absence of an
exact match as a limitation of the search performed, not as a claim of
priority.

We emphasize at the outset that the change in polynomial degree between
the first-order case (quadratic in $s$, affine in $w$) and the inertial
case (quartic in $s$, quadratic in $w$ on the reciprocal locus) is a
computed fact about each specific model, not an instance of a general
``dynamical-order law'' relating state dimension to reciprocal-polynomial
degree; no such general claim is made or needed here.

\section{Model}
\label{sec:model}

Consider two damped harmonic oscillators with masses $m_1,m_2$, natural
frequencies $\omega_1,\omega_2$ (in the original, mass-bearing
coordinates), raw damping constants $\gamma_1,\gamma_2$, and a linear
spring coupling of strength $k$. Each oscillator's momentum coordinate is
driven by an independent white-noise bath. The full state is
$(x_1,p_1,x_2,p_2)\in\R^4$.

\paragraph{Mass-weighted reduction.} The canonical transformation
$(x_i,p_i)\mapsto(\sqrt{m_i}\,x_i,\,p_i/\sqrt{m_i})$ has Jacobian
determinant exactly $1$, and reduces the general-mass problem to an
equivalent \emph{unit-mass} problem depending on the masses and raw
parameters only through the \emph{specific damping} rates
\[
\delta_i := \gamma_i/m_i,
\]
together with reduced frequencies absorbing the $1/m_i$ factors. Because
the transformation has unit Jacobian, it preserves every covariance
determinant appearing below; the palindromicity criterion we prove is
therefore exactly a statement about $\delta_1,\delta_2$, not about
$\gamma_i$ or $m_i$ individually. From here on, $\omega_1,\omega_2,k$
denote the (already reduced) unit-mass parameters, and we work entirely
in the unit-mass system; the final theorem is restated in terms of
$\gamma_i/m_i$ in Section~\ref{sec:restore}.

\paragraph{Drift and diffusion.} In the unit-mass system, with state
ordering $(x_1,p_1,x_2,p_2)$,
\[
A =
\begin{pmatrix}
0 & 1 & 0 & 0\\
-(\omega_1^2+k) & -\delta_1 & k & 0\\
0 & 0 & 0 & 1\\
k & 0 & -(\omega_2^2+k) & -\delta_2
\end{pmatrix},
\qquad
Q_1 = \operatorname{diag}(0,2\delta_1,0,0),
\qquad
Q_2 = \operatorname{diag}(0,0,0,2\delta_2).
\]
$Q_1,Q_2$ are the diffusion matrices for bath 1 and bath 2 acting alone
(unit temperature; a bath at temperature $T_i$ scales $Q_i$ by $T_i$, and
the full joint stationary covariance under both baths simultaneously,
$\Sigma(T_1,T_2)$, is linear in $T_1,T_2$ -- see
Section~\ref{sec:pencil}). The stationary covariances $C_1,C_2$ of the
two single-bath problems solve the Lyapunov equations
\begin{equation}
\label{eq:lyap}
A C_i + C_i A^{\mathsf T} = -Q_i, \qquad i=1,2.
\end{equation}

\paragraph{Relation to prior models.} Two damped harmonic oscillators,
each carrying its own position and momentum coordinate and each driven by
an independent heat bath, is not a new model class: the same $4\times4$
Lyapunov-equation setup (uniform damping, in that case) was used
by Boffi and De Gregorio~\cite{boffi2024} to study a resonance phenomenon
in the covariance entries as the ratio of the two spring constants
varies. Their analysis fixes a common damping rate throughout and does
not address palindromicity, reciprocal spectra, or any damping-ratio
criterion; the present work instead varies the two specific damping
rates $\delta_1,\delta_2$ independently and asks exactly when the
resulting covariance-pencil determinant is palindromic.

\paragraph{Admissible domain.} Throughout, $\omega_1,\omega_2,k,\delta_1,
\delta_2>0$ (open, all strict). The sole boundary degeneracy is $k=0$: at
$k=0$ the drift matrix $A$ becomes exactly block-diagonal, bath~1 cannot
influence oscillator~2 (and vice versa), and $C_1,C_2$ become rank-deficient
(singular). For $k\neq0$, the single-bath Kalman/Krylov controllability
determinant for each channel equals $k^2$ exactly, so $C_1,C_2\succ0$ for
every $k>0$ (Lemma~\ref{lem:pencil}). No further restriction on any
parameter is used anywhere below.

\section{Bath-response covariances and the pencil}
\label{sec:pencil}

\begin{lemma}[Pencil and positivity]
\label{lem:pencil}
With $Q(T_1,T_2):=T_1Q_1+T_2Q_2$ and $\Sigma(T_1,T_2)$ the stationary
solution of $A\Sigma+\Sigma A^{\mathsf T}=-Q(T_1,T_2)$,
\[
\Sigma(T_1,T_2) = T_1 C_1 + T_2 C_2.
\]
Writing $s:=T_1/T_2$ and $P(s):=\det(sC_1+C_2)$,
\[
P(s) = c_4 s^4+c_3s^3+c_2s^2+c_1s+c_0, \qquad c_4=\det C_1,\ \ c_0=\det C_2.
\]
For $k>0$, $C_1\succ0$ and $C_2\succ0$.
\end{lemma}

\begin{proof}[Proof sketch]
Linearity of $\Sigma$ in $(T_1,T_2)$ follows from linearity of the
Lyapunov equation~\eqref{eq:lyap} in the forcing term. $c_4,c_0$ are the
$s^4,s^0$ coefficients of $\det(sC_1+C_2)$, which equal $\det C_1,\det
C_2$ by multilinearity of the determinant in the pencil parameter.
Positive definiteness for $k>0$ follows from the exact controllability
determinant $k^2$ (single-bath Kalman rank criterion applied
to~\eqref{eq:lyap}) together with $A$'s stability for every
$\delta_1,\delta_2>0$: a stable, controllable pair yields a positive
definite stationary Gramian, a classical fact of linear systems theory
\cite{kailath1980}.
\end{proof}

\section{Palindromic sufficiency}
\label{sec:sufficiency}

\begin{lemma}[Sufficiency]
\label{lem:sufficiency}
If $\delta_1=\delta_2$ then $P(s)$ is palindromic: $c_4=c_0$ and
$c_3=c_1$, for every admissible $\omega_1,\omega_2,k>0$.
\end{lemma}

The shortest available proof substitutes $\delta_1=\delta_2=\delta$
directly into~\eqref{eq:lyap} before solving. This is a genuine
structural simplification, not merely a reduction in the number of free
symbols: with common specific damping, the $p$-block of $A$'s damping is
proportional to the identity, hence commutes with any orthogonal
transform diagonalizing the stiffness block
$\begin{psmallmatrix}\omega_1^2+k&-k\\-k&\omega_2^2+k\end{psmallmatrix}$.
This decouples the deterministic drift into two independent normal modes
with frequencies $\Omega_1,\Omega_2$ (the stiffness block's eigenvalues)
and common damping $\delta$, while the per-bath noise becomes a rank-one
source correlated across the two modes -- the same structure used for
Proposition~\ref{prop:slope} below. Solving~\eqref{eq:lyap} in these
coordinates gives $C_1(\delta),C_2(\delta)$ in closed form, and
$c_4-c_0=0$, $c_3-c_1=0$ are then confirmed by exact symbolic
simplification (Section~\ref{sec:cas}, item~1).

\section{The log-determinant comparator}

\begin{definition}
\label{def:F}
$F(\delta_1,\delta_2;\omega_1,\omega_2,k) := \log(\det C_2/\det C_1)$,
well-defined on the admissible interior by Lemma~\ref{lem:pencil}.
\end{definition}

$F=0$ is exactly bath-response volume balance ($\det C_1=\det C_2$), and
$\operatorname{sign}F$ distinguishes the two possible orderings of the
two single-bath covariance volumes. By Lemma~\ref{lem:sufficiency},
$F(\delta_2,\delta_2)=0$ for every admissible $\omega_1,\omega_2,k,
\delta_2$ -- the base point used throughout Sections~\ref{sec:slope}--\ref{sec:oriented}.

\section{Universal balance slope}
\label{sec:slope}

\begin{proposition}[Balance-locus slope]
\label{prop:slope}
At $\delta_1=\delta_2=\delta$,
\[
\left.\dd{F}{\delta_1}\right|_{\delta_1=\delta_2=\delta} = -\frac{4}{\delta},
\]
for every admissible $\omega_1,\omega_2,k,\delta>0$ -- independent of the
oscillator frequencies and the coupling strength.
\end{proposition}

\begin{proof}[Proof sketch]
Using the normal-mode decoupling of Section~\ref{sec:sufficiency}, the
adjoint-Lyapunov sensitivity identity gives, for $X_i:=\partial
C_i/\partial\delta_1$,
\[
\dd{F}{\delta_1} = \operatorname{tr}(C_2^{-1}X_2)-\operatorname{tr}(C_1^{-1}X_1)
 = 2\left[e_2^{\mathsf T}Y_2v_2 - e_2^{\mathsf T}Y_1u_1\right],
\]
where $e_2=(0,1,0,0)^{\mathsf T}$, $u_1=C_1e_2-e_2$, $v_2=C_2e_2$, and
$Y_i$ solves $A^{\mathsf T}Y_i+Y_iA=C_i^{-1}$ (this sign -- positive, not
negative -- was fixed after an initial sign error was caught by
cross-checking against direct differentiation; see
Section~\ref{sec:cas}). Substituting the closed-form $C_1(\delta),
C_2(\delta)$ from Section~\ref{sec:sufficiency} and solving for $Y_1,Y_2$
symbolically (fully general in $\omega_1,\omega_2,k,\delta$) yields the
stated closed form after exact simplification. The
$\omega_1,\omega_2,k$-cancellation arises from a combination of (i) the
equipartition structure of each normal mode's self-covariance, (ii) the
two channels sharing the identical denominator
$k^4/(\Delta D_{12}^2)$ at $\delta_1=\delta_2$ (with $\Delta,D_{12}$ as
in Section~\ref{sec:factorization}), which cancels in the derivative
ratio, and (iii) the shared specific damping making the rank-one
sensitivity $A'=-e_2e_2^{\mathsf T}$ act identically on both channels'
drift operator. This is a proved structural fact about this specific
model class, not claimed to generalize beyond it.
\end{proof}

\section{Global monotonicity}
\label{sec:factorization}

\begin{lemma}[Derivative factorization]
\label{lem:factorization}
\[
\dd{F}{\delta_1} = \frac{K_{\mathrm{total}}(\omega_1,\omega_2,k)\cdot
\delta_1^3\delta_2^4}{K(\delta_1,\delta_2)^4\cdot\det C_1\cdot\det C_2},
\]
where
\[
K(\delta_1,\delta_2) = a_1\delta_1^3\delta_2+a_2\delta_1^2\delta_2^2
+a_3\delta_1^2+a_4\delta_1\delta_2^3+a_5\delta_1\delta_2+a_6\delta_2^2,
\]
\[
a_1=\omega_2^2+k,\quad a_2=\omega_1^2+\omega_2^2+2k,\quad a_3=k^2,\quad
a_4=\omega_1^2+k,\quad a_5=(\omega_1^2-\omega_2^2)^2+2k^2,\quad a_6=k^2,
\]
and
\[
K_{\mathrm{total}}(\omega_1,\omega_2,k) = -\frac{4k^8}{\Delta^2},
\qquad
\Delta := k\omega_1^2+k\omega_2^2+\omega_1^2\omega_2^2.
\]
\end{lemma}

\begin{proof}[Proof sketch]
The six-monomial support and coefficient functions $a_1,\dots,a_6$ were
obtained by exact-rational interpolation across many
$(\omega_1,\omega_2,k)$ configurations and then proved as an identity by
exact zero-remainder polynomial division at independent, freshly chosen
points (Section~\ref{sec:cas}, item~2). $K_{\mathrm{total}}$ is
determined by evaluating the identity at the (fully general, symbolic)
balance locus $\delta_1=\delta_2=\delta$, where $K(\delta,\delta)=
\delta^2D_{12}$ with $D_{12}=(\Omega_1^2-\Omega_2^2)^2+2\delta^2
(\Omega_1^2+\Omega_2^2)$ the normal-mode cross-Sylvester-resolvent
denominator of Section~\ref{sec:sufficiency}, $\Delta=\Omega_1^2\Omega_2^2$
the product of squared normal-mode frequencies, and $\det C_1=\det C_2=
k^4/(\Delta D_{12}^2)$ there; combined with Proposition~\ref{prop:slope}'s
closed form $F'=-4/\delta$, this determines $K_{\mathrm{total}}$
uniquely (it is independent of $\delta$ by construction, so any one
evaluation point suffices), and the result is independently re-confirmed
by exact zero-remainder division against the general (non-symmetric)
identity at several fresh points, including off-diagonal ones
(Section~\ref{sec:cas}, item~3).
\end{proof}

\begin{theorem}[Global monotonicity and oriented comparator]
\label{thm:monotone}
Each $a_1,\dots,a_6>0$ by inspection (sums of squares and manifestly
positive terms; $a_5>0$ even at $\omega_1=\omega_2$, since $k>0$), so
$K>0$ for all $\delta_1,\delta_2>0$ (a sum of six positive terms).
$\Delta>0$ (a sum of three positive terms), so $K_{\mathrm{total}}<0$.
Since $\det C_1,\det C_2>0$ (Lemma~\ref{lem:pencil}), every factor in
Lemma~\ref{lem:factorization}'s identity has manifest sign, giving
\[
\dd{F}{\delta_1} < 0
\]
for every admissible interior point (fixed $\omega_1,\omega_2,k,\delta_2$).
Combined with the base point $F(\delta_2,\delta_2)=0$
(Section~\ref{sec:sufficiency}), $F$ is strictly decreasing in $\delta_1$
with a single zero crossing at $\delta_1=\delta_2$:
\[
\operatorname{sign}F = -\operatorname{sign}(\delta_1-\delta_2).
\]
Monotonicity of Gramian-type quantities under parameter perturbations is
a recurring theme in linear systems and sensitivity theory; we did not
locate a general theorem in the literature searched that directly implies
this specific determinant-ratio monotonicity for the pencil $sC_1+C_2$,
so we record it here as established for this model by the explicit
factorization of Lemma~\ref{lem:factorization} rather than as an instance
of a known general result.
\end{theorem}

\section{Oriented volume theorem}
\label{sec:oriented}

\begin{corollary}[Unique volume balance]
\label{cor:balance}
\[
\det C_1 = \det C_2 \iff \delta_1=\delta_2.
\]
\end{corollary}

This is immediate from Theorem~\ref{thm:monotone} and Definition~\ref{def:F}
(($F=0\iff\det C_1=\det C_2$)), and is a central result in its own right,
not merely a device toward palindromicity: it establishes that the
bath-response covariance \emph{volumes} $\det C_1,\det C_2$ coincide
exactly once, at exactly the equal-specific-damping locus, with a definite
orientation on either side.

\section{Global palindromicity theorem}
\label{sec:main}

\begin{theorem}[Palindromicity iff equal specific damping]
\label{thm:main}
For the admissible genuinely coupled two-mode inertial OU system,
\[
P(s) \text{ is palindromic} \iff \delta_1=\delta_2.
\]
\end{theorem}

\begin{proof}
\emph{Necessity.} If $P(s)$ is palindromic then $c_4=c_0$, i.e.\ $\det
C_1=\det C_2$ (Lemma~\ref{lem:pencil}), hence $\delta_1=\delta_2$
(Corollary~\ref{cor:balance}). The coefficient identity $c_3=c_1$ is not
invoked. \emph{Sufficiency.} Lemma~\ref{lem:sufficiency}.
\end{proof}

\subsection{Restoring general masses}
\label{sec:restore}

Restoring $\delta_i=\gamma_i/m_i$ (Section~\ref{sec:model}; the
mass-weighted reduction is exact and preserves every determinant used
above):
\[
P(s)\text{ is palindromic} \iff \frac{\gamma_1}{m_1}=\frac{\gamma_2}{m_2}.
\]

\section{Spectral interpretation}

\begin{corollary}[Reciprocal generalized eigenvalues]
For $k>0$, $N:=C_1^{-1}C_2$ is similar to a symmetric positive-definite
matrix (via $C_1^{-1/2}$), hence has four strictly positive real
eigenvalues $\mu_1,\dots,\mu_4$ -- the roots of $P(s)=0$ are $s=-\mu_i$.
$P(s)$ is palindromic if and only if $\{\mu_1,\dots,\mu_4\}$ splits into
two reciprocal pairs $\{\mu,1/\mu\},\{\nu,1/\nu\}$. Combined with
Theorem~\ref{thm:main}:
\[
\delta_1=\delta_2 \iff P(s)\text{ palindromic} \iff \{\mu_i\}\text{ splits into reciprocal pairs.}
\]
\end{corollary}
We emphasize the convention: eigenvalues are those of $C_1^{-1}C_2$
(equivalently, of the symmetrized matrix), not $C_2^{-1}C_1$; the sign
convention $s=-\mu_i$ in the palindromic-root characterization follows
from writing $P(s)=\det C_1\cdot\det(sI+C_1^{-1}C_2)$. Reciprocal
eigenvalue pairing for structurally palindromic matrix pencils (those
built from an explicit relation such as $B=A^{\mathsf T}$) is classical
matrix-pencil theory~\cite{mackey2006}; the pencil $sC_1+C_2$ studied
here carries no such a priori structural relation between $C_1$ and
$C_2$ (each solves an independent Lyapunov equation), so the reciprocal
pairing recorded in this corollary is a \emph{derived} consequence of
Theorem~\ref{thm:main}, not an instance of the structural theory.

\section{Reciprocal coordinate}

On the palindromic locus ($\delta_1=\delta_2$), dividing $P(s)$ by $s^2$
and using the classical reciprocal-equation identity $s^2+s^{-2}=w^2-2$
with $w:=s+s^{-1}$,
\[
P(s)/s^2 = c_4w^2+c_3w+(c_2-2c_4),
\]
a genuine (non-degenerate) quadratic in $w$, with $c_4,c_3,c_2$ explicit
rational functions of $\omega_1,\omega_2,k,\delta$. This is the inertial
analogue of the affine reciprocal-coordinate reduction in the first-order
OU problem~\cite{dejesus2026heterogeneity}; the change in degree (affine
versus quadratic) is a fact about the two specific models, not an
instance of a general law relating state dimension to reduced-polynomial
degree.

\section{Discussion}

\emph{Specific damping, not raw parameters, controls reciprocal
symmetry.} The palindromicity criterion is $\gamma_1/m_1=\gamma_2/m_2$,
not $\gamma_1=\gamma_2$ or $m_1=m_2$ individually -- a direct consequence
of the mass-weighted reduction (Section~\ref{sec:model}) having unit
Jacobian. \emph{Volume balance is unique} (Corollary~\ref{cor:balance}):
the two single-bath covariance determinants coincide at exactly one
locus. \emph{The comparator is oriented}: $F$ does more than vanish at
balance -- its sign tracks which bath produces the larger stationary
covariance volume, throughout the whole domain
(Theorem~\ref{thm:monotone}). \emph{The balance crossing has a universal
slope}, $-4/\delta$, independent of frequencies and coupling
(Proposition~\ref{prop:slope}). \emph{Reciprocal spectral pairing}
(Section~7) gives a coordinate-independent restatement of the same
criterion. The decomposition of the monotonicity identity into canonical
internal response terms ($F'=2(R_2-R_1)$ in the notation of
Proposition~\ref{prop:slope}'s proof) may provide a direction for future
work.

\subsection{Relation to the first-order theory}

The first-order OU reciprocity result~\cite{dejesus2026heterogeneity} (see
also this repository's Lean formalization) concerns a $2\times2$
stationary covariance, a quadratic pencil determinant, a trace-normalized
forcing-ratio criterion, and an affine reduction in $w$. The inertial
system studied here has a $4\times4$ covariance, a quartic pencil
determinant, an equal-specific-damping criterion, a quadratic
$w$-reduction on the reciprocal locus, and the additional oriented-volume
theorem with no first-order analogue established here. The inertial
result is not a routine extension of the first-order one: the
palindromicity criterion, the reduced polynomial's degree, and the proof
architecture (normal-mode decoupling, adjoint-Lyapunov sensitivity) are
each specific to the inertial state-space structure.

\subsection{Limitations}

This work concerns a linear Gaussian system with exactly two inertial
modes, stationary statistics only, and the specific per-bath forcing
structure of Section~\ref{sec:model} (each bath coupled to exactly one
oscillator's momentum coordinate). Results are established on the
genuinely coupled interior $k>0$; no claim is made at $k=0$. No claim is
made about nonlinear extensions, about networks of more than two
oscillators, or that the exact slope $-4/\delta$ holds outside this
specific model class.

\section{Computer-assisted proof disclosure}
\label{sec:cas}

Several polynomial identities used above were verified by exact computer
algebra. These checks use symbolic or exact-rational arithmetic rather
than floating-point sampling; the accompanying verification scripts
(\texttt{verification/inertial/} in the project repository) reproduce
every load-bearing identity from a single documented entry point,
\texttt{verify\_global\_chain.py}.

\begin{enumerate}
\item \textbf{Sufficiency} (Lemma~\ref{lem:sufficiency}): exact symbolic
computation -- $C_1(\delta),C_2(\delta)$ solved via exact linear algebra
with $\delta_1=\delta_2=\delta$ substituted before solving;
$c_4-c_0,c_3-c_1$ confirmed to simplify symbolically to $0$.
(\texttt{verify\_sufficiency.py})
\item \textbf{The six-coefficient $K$ identity} (Lemma~\ref{lem:factorization}):
discovered via exact-rational interpolation, proved by exact
zero-remainder polynomial division against the independently computed
true numerator, at fresh points not used in the interpolation.
\item \textbf{$K_{\mathrm{total}}=-4k^8/\Delta^2$} (Lemma~\ref{lem:factorization}):
derived from the balance-slope proposition, independently re-confirmed
by exact zero-remainder division at fresh points including an
off-diagonal ($\delta_1\neq\delta_2$) end-to-end check.
(\texttt{verify\_general\_Fprime.py})
\item \textbf{The balance-slope proposition} (Proposition~\ref{prop:slope}):
fully symbolic exact derivation (not sampled) via the adjoint-Lyapunov
route, with the sign convention on the $Y_i$ equation independently
cross-checked against direct symbolic differentiation.
(\texttt{verify\_balance\_slope.py})
\end{enumerate}

Given these closed forms, $K>0$, $K_{\mathrm{total}}<0$, $F'<0$, the
oriented sign theorem, and palindromicity necessity are each short,
human-checkable steps (Section~\ref{sec:factorization}--\ref{sec:main}).
Earlier stages of this investigation used large parameter sweeps (a
729-point rational search, 215 and 102 exact line searches, 17
six-monomial discovery configurations, 66 scalar-sign spot checks) to
locate the closed forms above before they were proved in general; that
search history carries no theorem weight and is not part of this proof.

\section{Reproducibility}

All symbolic computation uses SymPy with exact (arbitrary-precision
integer and rational) arithmetic; no floating-point computation is used
in any load-bearing step. The verification workflow is:
\begin{center}
\texttt{python verification/inertial/verify\_global\_chain.py}
\end{center}
which reproduces, from a single command, every identity disclosed in
Section~\ref{sec:cas} together with a fresh end-to-end regression check
(new parameter points, not used in discovery) confirming
positive-definiteness, the base point, global monotonicity, the oriented
sign theorem, and palindromicity exactly at $\delta_1=\delta_2$ together
with its exact failure at representative $\delta_1\neq\delta_2$ points.
Proof-generating and verification scripts (\texttt{verification/inertial/})
are kept separate from this project's historical exploratory search
scripts, which are not included in this repository.

\section*{AI assistance disclosure}

The research and conceptual development are by Elias De Jes\'us, with AI
assistance used in symbolic-workflow development, code generation and
debugging, adversarial proof auditing, and manuscript editing.
Mathematical claims in this manuscript were checked through the
documented exact computational workflow of Section~\ref{sec:cas} and
Section~9; independent human derivation is not claimed for the
identities disclosed there as computer-assisted.

\section*{Acknowledgments}

None.

\bibliographystyle{plain}
\bibliography{references}

\end{document}
